\newif\ifglobal

%%%%%%%%%%%%%%%%%%%%%%%
%%% Compile to view %%%
%%%%%%%%%%%%%%%%%%%%%%%

%%%%%%%%%%%%%%%%%%%%%%%%%%%%%%%%%%%%%%%%%%%%%%%%%%%
%%% Readme:                                     %%%
%%% https://www.overleaf.com/read/bwdjvfjksvjy/ %%%
%%%%%%%%%%%%%%%%%%%%%%%%%%%%%%%%%%%%%%%%%%%%%%%%%%%

% >>> M?: marks a module setting number or important notice


% >>> M4: Set {./relative-path} to external files for this module <<<
\def \modulefiles {./40-H264ENC}
% To add an external file, use:
% (Replace '---' with actual filename)
% '\input{\modulefiles/tables/---.tex}' for tables
% '\includegraphics{\modulefiles/figures/---}' for figures
% '\subfile{\modulefiles/---}' for register files

% >>> M1: This module is in global project? <<<
% 'YES' - uncomment; 'NO' - comment out
\globaltrue

\ifglobal
    \documentclass[main\_\_CN.tex]{subfiles}

\else
    % Font "FandolSong-Regular" used by ctex does not have GSUB table.
    % It causes warnings that do not affect compilation and can be ignored.
    \PassOptionsToPackage{quiet}{fontspec}

    \documentclass[openany, 10pt]{book}

    % >>> M2: In ./00-shared/preamble-trm-module, also find
    %         settings for visibility of todo notes, watermark <<<
    \usepackage{preamble-trm-module}


    % Enable Chinese typesetting
    \usepackage{ctex}

    % >>> M3: Temporary labels in Chinese
    %         you can add your own labels there <<<
    \myexternaldocument{./00-shared/temp-labels-cn}

\fi %\ifglobal


\setcounter{secnumdepth}{4}

% >>> M5: If this module needs any updates in
% './00-shared/preamble-trm-module.sty'
% (include additional package, etc.), contact your module PM
% or create the file './00-shared/changelog.tex' make the
% necessary updates yourself and document them in this file <<<

% Make text left-aligned, comment out for full justification
\raggedright


\begin{document}


% Sets language into which pre-defined bits of text are translated
\selectlanguage{chinese}

% Do not delete this block
\ifglobal
% Add the following front matter if
% this module is outside of global project
\else
    \InputIfFileExists{readme}{}{}
    {\let\clearpage\relax \listoftodos}
    {\let\clearpage\relax \listofdonetodos}
    \tableofcontents
    %\listoftables
    %\listoffigures
\fi


%%%%%%%%%%%%%%%%%%%%%%%%%
%%% TRM Module Begins %%%
%%%%%%%%%%%%%%%%%%%%%%%%%

\hypertarget{h264enc}{}
\chapter{H264 编码器}\label{mod:h264enc}

\section{概述}

H264 编码器基于 H264 标准，用于视频序列的实时压缩。使用 H264 编码器可以显著减少数据总量，同时尽量减少视频质量的损失。



\section{术语}

\begin{longtable}[c]{ R{4cm} p{12cm} }
    \textbf{比特流} & 表示一个或多个编码视频序列的比特序列。 \\
    \textbf{配置文件} & 由 H264 标准规定的整个比特流语法的子集。 \\
    \textbf{级别} & 在每个配置文件中指定，用于约束比特流中语法元素的值。 \\
    \textbf{宏块} & 即 Macro Block，简称 MB。对于一个具有三个样本阵列的图像，其 MB 由一个 16 x 16 像素的亮度样本块和两个对应的色度样本块组成。 \\
    \textbf{切片} & 每个切片包含整数个宏块。 \\
    \textbf{I 帧} & 采用帧内编码的帧。 \\
    \textbf{P 帧} & 采用帧间预测编码的帧。 \\
    \textbf{I 切片} & I 帧中的切片。 \\
    \textbf{P 切片} & P 帧中的切片。 \\
    \textbf{P-skip 宏块} & 在 P 帧中跳过编码的宏块。 \\
    \textbf{CAVLC} & 基于上下文的自适应可变长度编码。 \\
    \textbf{MV} & 运动矢量。 \\
    \textbf{GOP} & 图像组。在一个 GOP 中，第一帧为 I 帧，后续帧为 P 帧。 \\
    \textbf{QP} & 量化参数。 \\
    \textbf{ROI} & 感兴趣区域。 \\
    \textbf{AXI} & 高级可扩展接口。 \\
    \textbf{TX} & 发送。 \\
    \textbf{RX} & 接收。 \\
    \textbf{DB} & 去块滤波器。\\
    \textbf{NUM\_H\_MB} & 图像中水平宏块的数量。\\
    \end{longtable}

%%%%%%%%%%%%%%%%
%%% Features %%%
%%%%%%%%%%%%%%%%

\section{特性}
\label{sec:h264enc-features}

H264 编码器具有如下特性：
\begin{itemize}
    \item 基于基线配置文件，支持级别 4.1。
    \item 支持 8 位 YUV 4:2:0 逐行视频，最大分辨率达 1920 × 1088 @ 30 fps。
    \item 支持 I 帧和 P 帧编码。
    \item 提供 GOP 和双流两种工作模式。在双流模式下，两个待编码的视频序列的总带宽不应超过 1920 x 1088 @ 30fps。
    \item 内部帧支持 4 x 4 和 16 x 16 的亮度宏块分割。
    \item 帧内预测中，4 x 4 的亮度宏块支持 9 种预测模式，16 x 16 的亮度宏块支持 4 种预测模式。
    \item 帧内预测中，色度宏块支持 4 种预测模式。
    \item 帧间预测中，亮度宏块支持多种分割大小：4 x 4、4 x 8、8 x 4、8 x 8、8 x 16、16 x 8、16 x 16。
    \item 运动估算精度达 1/2 和 1/4 像素。
    \item 水平运动搜索范围为 [-29.75, +16.75]，垂直运动搜索范围为 [-13.75, +13.75]。
    \item 配备去块滤波器。
    \item 使用 CAVLC 编码。
    \item 支持 P-skip 宏块。
    \item 支持在 P 片中包含 I 宏块。
    \item 支持亮度和色度分量量化结果降采样操作。
    \item 可配置固定 QP，并支持宏块级别的码率控制。
    \item 提供 MV 合并功能。当宏块的 MV 大于 0 时，此 MV 会输出到内存中。
    \item 支持 ROI 功能，最多可配置 8 个矩形 ROI，位置任意（允许重叠，优先级固定）。每个 ROI 可配置一个常量 QP 或 QP 偏移，非 ROI 区域可以配置 QP 偏移。
\end{itemize}


%%%%%%%%%%%%%%%%%%%%
%%% Architecture %%%
%%%%%%%%%%%%%%%%%%%%

\section{架构}
\label{sec:h264enc-arch-overview}

H264 编码器的架构如图 \ref{fig:h264-architecture} 所示。

\begin{figure}[H]
    \centering
    \includegraphics[width=0.8\textwidth]{\modulefiles/figures/h264-arc.eps}
    \caption{H264 编码器架构}
    \label{fig:h264-architecture}
\end{figure}

H264 编码器由两部分组成：H264 专用 DMA (H264\_DMA) 和编码器算法核 (ENC\_CORE)。H264\_DMA 用于数据传输，通过两组 AXI 总线与内部存储器和外部存储器连接，同时通过多个 DMA 发送/接收通道与 ENC\_CORE 进行通信。具体而言，H264\_DMA 配备 5 个 TX 通道 \regindex{n} 和 6 个 RX 通道 \regindex{n}，其中 TX 通道 3、TX 通道 4 和 RX 通道 5 为 H264\_DMA 的内部通道，不与 ENC\_CORE 连接。ENC\_CORE 用于视频序列的压缩处理。

H264\_DMA 的各个发送通道和接收通道功能如下：
\begin{itemize}
\item \textbf{TX 通道 0} 将原始图像发送到 ENC\_CORE。
\item \textbf{TX 通道 1} 将参考图像发送到 ENC\_CORE。
\item \textbf{TX 通道 2} 将去块滤波器中间数据发送到 ENC\_CORE。
\item \textbf{TX 通道 3} 和 \textbf{TX 通道 4} 将参考图像（即上一次去块滤波后的图像）发送到 \textbf{RX 通道 5}。
\item \textbf{RX 通道 0} 和 \textbf{RX 通道 1} 将去块滤波器的结果（即参考图像）发送到外部存储器。
\item \textbf{RX 通道 2} 将去块滤波器中间数据发送到内部存储器。
\item \textbf{RX 通道 3} 将 MV 合并结果发送到内部存储器。
\item \textbf{RX 通道 4} 将编码的比特流发送到外部存储器。
\item \textbf{RX 通道 5} 将参考图像（即上一次去块滤波后的图像）发送到内部存储器。
\end{itemize}


根据描述可知，H264\_DMA 通道可分为两类：

\begin{itemize}
\item 访问内部存储（通过 AXI\_M1 总线）：TX 通道 1/2，RX 通道 2/3/5
\item 访问外部存储（通过 AXI\_M2 总线）：TX 通道 0/3/4，RX 通道 0/1/4
\end{itemize}

%%%%%%%%%%%%%%%%%%%%%%%%%%%%%%
%%% Functional Description %%%
%%%%%%%%%%%%%%%%%%%%%%%%%%%%%%

\section{功能描述}
\label{sec:h264enc-func-descr}

\subsection{编码器算法核 (ENC\_CORE)}\label{sec:h264-core}
编码器算法核根据软件配置对视频序列进行压缩和编码。


\subsubsection{架构}
编码器算法核的架构如图 \ref{fig:h264-enc-core} \textit{\nameref{fig:h264-enc-core}} 所示。

\begin{figure}[H]
    \centering
    \includegraphics[width=1.0\textwidth]{\modulefiles/figures/h264-enc-core}
    \caption{ENC\_CORE 架构}
    \label{fig:h264-enc-core}
\end{figure}

每个模块的功能如下：
\begin{itemize}
    \item register：选择和配置寄存器。
    \item top\_ctrl：控制和调度编码器算法核的每个模块。
    \item cur\_mb：从 H264\_DMA 的 TX 通道 0 (tx\_pop\_ori) 中读取原始图像宏块，并分发给 intra、ime、fme 和 mc 模块。
    \item intra：使用 trans\_quant\_decimate 模块计算的重建结果对宏块进行帧内预测，将帧内预测的结果和残差发送给 trans\_quant\_decimate 模块，并将帧内预测成本发送给 rdo 模块。
    \item fetch：从 H264\_DMA 的 TX 通道 1 (tx\_pop\_ref) 中读取搜索范围内的参考图像宏块，并将其分发给 ime、fme 和 mc 模块。
    \item ime：接收来自 cur\_mb 模块的原始图像宏块，接收来自 fetch 模块的参考图像宏块，在搜索范围内进行整数运动估计，并将整数运动估计结果发送到 fme 模块。
    \item fme：接收来自 cur\_mb 模块的原始图像宏块，接收来自 fetch 模块的参考图像宏块，接收来自 ime 模块的整数运动估计结果，并在整数运动估计的基础上进行分数运动估计。将分数运动估计结果发送到 mc 模块，并将帧间预测成本发送给 rdo 模块。
    \item mc：接收来自 cur\_mb 模块的原始图像宏块，接收来自 fetch 模块的参考图像宏块，接收来自 fme 模块的分数运动估计结果，根据分数运动估计结果进行运动补偿。将帧间预测结果和残差之和发送到 trans\_quant\_decimate 模块，并将运动矢量发送到 mv\_merge 模块。
    \item rdo：接收来自 intra 模块的帧内预测成本，接收来自 fme 模块的帧间预测成本，选择当前宏块的最终类型，支持 P 切片中的 I 宏块功能，并将最终的宏块类型发送到 trans\_quant\_decimate 模块。
    \item trans\_quant\_decimate：先接收来自 rdo 模块的最终宏块类型，接收来自 intra 模块的帧内预测和残差结果，接收来自 mc 模块的帧间预测和残差结果，然后进行 DCT、量化、反量化、反 DCT 和量化结果降采样，最后将残差信息发送到 rate\_ctrl 模块，将降采样的量化结果发送到 cavlc 模块，并将重建结果发送到 deblock 模块。
    \item cavlc：接收来自 trans\_quant\_decimate 模块的量化结果和当前宏块的预测信息，进行上下文自适应变长编码（CAVLC），并将编码流写入 H264\_DMA 的 RX 通道 4 (rx\_push\_bs) 中。
    \item deblock：接收来自 trans\_quant\_decimate 模块的重建结果和当前宏块的预测信息，进行去块滤波，并将去块滤波结果（即参考图像）写入 H264\_DMA 的 RX 通道 0 (rx\_push\_db\_12line) 和 RX 通道 1 (rx\_push\_db\_4line) 中。此模块通过 H264\_DMA 的 TX 通道 2（tx\_pop\_db\_tmp）和 RX 通道 2 通道（rx\_push\_db\_tmp）读写去块滤波的中间数据。
    \item mv\_merge：接收来自 mc 模块的运动矢量 MV，进行 MV 合并，并将合并结果（宏块运动信息）写入 H264\_DMA 的 RX 通道 3 (rx\_push\_mv\_merge)。
    \item rate\_ctrl：接收来自 trans\_quant\_decimate 模块的残差信息，接收来自 cavlc 模块的编码信息，以及 register 模块配置的帧级信息，进行宏块级码率控制，并将计算出的量化参数 QP 传递给 roi 模块。
    \item roi：接收来自 rate\_ctrl 模块的量化参数，接收由 register 模块配置的 ROI 信息，调整 QP，并使用调整后的 QP 处理当前宏块。
\end{itemize}



\subsubsection{量化结果降采样}\label{sec:h264-decimate}
ENC\_CORE 根据每个宏块的最终量化结果的大小计算宏块的得分，然后将得分与可配置的阈值寄存器进行比较。如果低于阈值，量化结果将被设置为 0。否则，量化结果将保持不变，以进一步提高压缩比。对每个宏块的亮度和色度分量独立执行上述过程。不同类型的宏块有不同的比较方法，如下所示：

\begin{itemize}
    \item 亮度分量的 I 宏块，采用 16 x 16 分区：

    \begin{center}
    $\textrm{l\_score\_mb\_cmp} + \textrm{\hyperref[fielddesc:H264AI16X16DECISCOREOFFSET]{H264\_\regindex{x}\_I16X16\_DECI\_SCORE\_OFFSET}} < \textrm{\hyperref[fielddesc:H264ALDECISCORE]{H264\_\regindex{x}\_L\_DECI\_SCORE}}$
    \end{center}


    其中，l\_score\_mb\_cmp 是 ENC\_CORE 计算得出的当前宏块亮度分量的得分。 \hyperref[fielddesc:H264AI16X16DECISCOREOFFSET]{H264\_\regindex{x}\_I16X16\_DECI\_SCORE\_OFFSET} 为可配置的偏移寄存器。 \hyperref[fielddesc:H264ALDECISCORE]{H264\_\regindex{x}\_L\_DECI\_SCORE} 是可配置的阈值寄存器。\regindex{x} 可以是 A 或 B，表示对应于视频序列 A 或 B。当满足以上条件时，当前亮度分量宏块的量化结果将被降采样为 0。
    \item 亮度分量 I 宏块，4 x 4 分区：当前亮度分量宏块的量化结果不会被减少为 0。
    \item 色度分量 I 宏块：

    \begin{center}
    $\textrm{c\_score\_mb\_cmp} + \textrm{\hyperref[fielddesc:H264AICHROMADECISCOREOFFSET]{H264\_\regindex{x}\_I\_chroma\_DECI\_SCORE\_OFFSET}} < \textrm{\hyperref[fielddesc:H264ACDECISCORE]{H264\_\regindex{x}\_C\_DECI\_SCORE}}$
    \end{center}

    其中，c\_score\_mb\_cmp 是 ENC\_CORE 计算的当前色度分量宏块的得分。 \hyperref[fielddesc:H264AICHROMADECISCOREOFFSET]{H264\_\regindex{x}\_I\_chroma\_DECI\_SCORE\_OFFSET} 是可配置的偏移寄存器。 \hyperref[fielddesc:H264ACDECISCORE]{H264\_\regindex{x}\_C\_DECI\_SCORE} 是可配置的阈值寄存器。\regindex{x} 可以是 A 或 B，表示对应于视频序列 A 或 B。当满足以上条件时，当前色度分量宏块的量化结果将被降采样为 0。
    \item 亮度分量 P 宏块:

    \begin{center}
    $\textrm{l\_score\_mb\_cmp} + \textrm{\hyperref[fielddesc:H264AP16X16DECISCOREOFFSET]{H264\_\regindex{x}\_P16X16\_DECI\_SCORE\_OFFSET}} < \textrm{\hyperref[fielddesc:H264ALDECISCORE]{H264\_\regindex{x}\_L\_DECI\_SCORE}}$
    \end{center}

    其中，l\_score\_mb\_cmp 是由 ENC\_CORE 计算的当前亮度分量宏块的得分。\hyperref[fielddesc:H264AP16X16DECISCOREOFFSET]{H264\_\regindex{x}\_P16X16\_DECI\_SCORE\_OFFSET} 是可配置的偏移寄存器。\hyperref[fielddesc:H264ALDECISCORE]{H264\_\regindex{x}\_L\_DECI\_SCORE} 是可配置的阈值寄存器。\regindex{x} 可以是 A 或 B，表示对应于视频序列 A 或 B。当满足以上条件时，当前亮度分量宏块的量化结果将被降采样为 0。
    \item 色度分量 P 宏块:

    \begin{center}
    $\textrm{c\_score\_mb\_cmp} + \textrm{\hyperref[fielddesc:H264APCHROMADECISCOREOFFSET]{H264\_\regindex{x}\_P\_chroma\_DECI\_SCORE\_OFFSET}} < \textrm{\hyperref[fielddesc:H264ACDECISCORE]{H264\_\regindex{x}\_C\_DECI\_SCORE}}$
    \end{center}

    其中，c\_score\_mb\_cmp 是由 ENC\_CORE 计算的当前色度分量宏块的得分。 \hyperref[fielddesc:H264APCHROMADECISCOREOFFSET]{H264\_\regindex{x}\_P\_chroma\_DECI\_SCORE\_OFFSET} 是可配置的偏移寄存器。 \hyperref[fielddesc:H264ACDECISCORE]{H264\_\regindex{x}\_C\_DECI\_SCORE} 是可配置的阈值寄存器。 \regindex{x} 可以是 A 或 B，表示对应于视频序列 A 或 B。当满足以上条件时，当前色度分量宏块的量化结果将被降采样为 0。
\end{itemize}


\subsubsection{运动矢量 (MV) 合并}\label{sec:h264-mv-merge}
ENC\_CORE 根据与 MV 合并相关的寄存器配置，合并亮度宏块的 MV，并将其写出到 H264\_DMA 的 RX 通道 3 中 (rx\_push\_mv\_merge)，最终存储在内部存储器中供上层应用使用，例如运动检测等。

与 MV 合并相关的寄存器包括：
\begin{itemize}
    \item \hyperref[fielddesc:H264AMVMERGEEN]{H264\_\regindex{x}\_MV\_MERGE\_EN}：此寄存器控制 MV 合并功能的启用。将此寄存器设置为 1 表示启用相应视频序列的 MV 合并功能。将此寄存器设置为 0 表示关闭相应视频序列的 MV 合并功能。\regindex{x} 可以是 A 或 B，表示对应于视频序列 A 或 B。即使启用了 MV 合并功能，I 宏块也不会向存储器输出任何数据，因为 I 宏块 没有 MV。
    \item \hyperref[fielddesc:H264MVMERGETYPE]{H264\_MV\_MERGE\_TYPE}：此寄存器控制 MV 合并的类型。将此寄存器配置为 0， 表示只保留最终类型为 P 宏块且亮度分量分区为 16 x 16 的宏块的 MV。将此寄存器配置为 1，表示只保留最终类型为 P 宏块的模最小的 MV。将此寄存器配置为 2， 表示只保留最终类型为 P 宏块的模最大的 MV。3 是无效值。当寄存器 \hyperref[fielddesc:H264AMVMERGEEN]{H264\_\regindex{x}\_MV\_MERGE\_EN} 设置为 1 时，此寄存器的配置生效。
    \item \hyperref[fielddesc:H264INTMVOUTEN]{H264\_INT\_MV\_OUT\_EN}：此寄存器控制最终存储在内存中的合并 MV 的类型。将此寄存器设置为 0，表示如果保留的合并 MV 不为 0 时将被存储在内存中，否则将被丢弃。将此寄存器设置为 1，表示如果保留的合并 MV 的整数部分不为 0 时将被存储在内存中，否则将被丢弃。当寄存器 \hyperref[fielddesc:H264AMVMERGEEN]{H264\_\regindex{x}\_MV\_MERGE\_EN} 设置为 1 时，此寄存器的配置生效。
    \item \hyperref[fielddesc:H264MBVALIDNUM]{H264\_MB\_VALID\_NUM}：此寄存器表示存储在内存中的当前图像 MV 合并结果的数量。一个宏块对应一个 MV 合并结果。当寄存器 \hyperref[fielddesc:H264AMVMERGEEN]{H264\_\regindex{x}\_MV\_MERGE\_EN} 设置为 1 时，此寄存器的值有意义，否则此寄存器的值无意义且保持为 0。
\end{itemize}


\subsubsection{宏块级码率控制}\label{sec:h264-rate-ctrl}
ENC\_CORE 根据码率控制相关的寄存器的配置、当前宏块的编码位数和残差的平均绝对值 (MAD)，计算出每个宏块的 QP 调整量。然后使用计算得出的 QP 调整量和寄存器中配置的帧级 QP，为每个宏块确定最终的 QP，从而实现编码位数的调整。


宏块级码率控制相关的寄存器包括：
\begin{itemize}
    \item \hyperref[fielddesc:H264AMBRATECTRLEN]{H264\_\regindex{x}\_MB\_RATE\_CTRL\_EN}：此寄存器控制是否启用宏块级码率控制。将此寄存器设置为 1，表示启用相应视频序列的宏块级码率控制。将此寄存器设置为 0，表示关闭相应视频序列的宏块级码率控制。\regindex{x} 可以是 A 或 B，表示对应于视频序列 A 或 B。注意，由于 I 帧使用恒定 QP 编码，因此即使启用了宏块级码率控制，I 帧中每个宏块的 QP 值也不会改变。
    \item \hyperref[fielddesc:H264ARATECTRLU]{H264\_\regindex{x}\_RATE\_CTRL\_U}：此寄存器的配置值由软件中的帧级码率控制计算得出。ENC\_CORE 需要此值来计算每个宏块的 QP 调整量。\regindex{x} 可以是 A 或 B，表示对应于视频序列 A 或 B。此寄存器配置值的计算方法属于软件中的帧级码率控制算法范畴，本章不再讨论。
    \item \hyperref[fielddesc:H264AMADFRAMEPRED]{H264\_\regindex{x}\_MAD\_FRAME\_PRED}：此寄存器是帧级 MAD 的预测值，由软件中的帧级码率控制进行配置。ENC\_CORE 利用此值来计算每个宏块的 QP 调整量。\regindex{x}可以是 A 或 B，表示对应于视频序列 A 或 B。此寄存器配置值的计算方法属于软件中的帧级码率控制算法的范畴，本章不再讨论。
    \item \hyperref[fielddesc:H264AQP]{H264\_\regindex{x}\_QP}：此寄存器是帧级 QP 值，值范围是 10 到 51，由软件中的帧级码率控制进行配置。ENC\_CORE 根据此值来调整每个宏块的 QP，从而实现宏块级别的码率控制。\regindex{x} 可以是 A 或 B，表示对应于视频序列 A 或 B。此寄存器的配置值的计算方法属于软件中的帧级码率控制算法的范畴，本章不再讨论。
    \item \hyperref[fielddesc:H264AQPMAX]{H264\_\regindex{x}\_QP\_MAX}：此寄存器用于限制宏块级码率控制最终计算的 QP 的最大值，值范围为 10 到 51。当宏块级码率控制最终计算的 QP 大于此寄存器配置值时，将被限制为此寄存器的配置值，否则保留原值。\regindex{x} 可以是 A 或 B，表示对应于视频序列 A 或 B。
    \item \hyperref[fielddesc:H264AQPMIN]{H264\_\regindex{x}\_QP\_MIN}：此寄存器用于限制宏块级码率控制最终计算的 QP 的最小值，值范围为 10 到 51。如果宏块级码率控制最终计算的 QP 小于此寄存器配置值，将被限制为此寄存器的配置值，否则保留原值。\regindex{x} 可以是 A 或 B，表示对应于视频序列 A 或 B。
    \item \hyperref[fielddesc:H264FRAMEMADSUM]{H264\_FRAME\_MAD\_SUM}：此寄存器表示当前帧中所有宏块的 MAD 累积和。软件需要在完成一帧视频的编码后读取此寄存器的值，进行帧级码率控制相关的计算。软件中的帧级码率控制算法超出了本章的范畴，不再讨论。
    \item \hyperref[fielddesc:H264FRAMEENCBITS]{H264\_FRAME\_ENC\_BITS}：此寄存器表示当前帧中所有宏块的编码位数的总和。软件需要在完成一帧视频的编码后读取此寄存器的值，进行帧级码率控制。软件中的帧级码率控制算法超出了本章的范畴，不再讨论。
    \item \hyperref[fielddesc:H264FRAMEQPSUM]{H264\_FRAME\_QP\_SUM}：此寄存器表示当前帧中所有宏块的 QP 总和。软件需要在完成一帧视频的编码后读取此寄存器的值，进行帧级码率控制。软件中的帧级码率控制算法超出了本章的范畴，不再讨论。
\end{itemize}


\subsubsection{感兴趣区域 (ROI)}\label{sec:h264-roi}
根据相关寄存器的配置，ENC\_CORE 在一帧视频中最多支持 8 个可重叠的固定优先级矩形 ROI。通过调整 QP（由宏块级码率控制计算得出），每个 ROI 的编码图像质量各不相同。

与 ROI 相关的寄存器如下：
\begin{itemize}
    \item \hyperref[fielddesc:H264AROIEN]{H264\_\regindex{x}\_ROI\_EN}：此寄存器用于开启或关闭 ROI 功能。将其设置为 1 以启用 ROI 功能，设置为 0 禁用 ROI 功能。\regindex{x} 可以是 A 或 B，表示对应于视频序列 A 或 B。
    \item \hyperref[fielddesc:H264AROIMODE]{H264\_\regindex{x}\_ROI\_MODE}：此寄存器用于指定 ROI 的工作模式。如果配置为 0，ROI 将以固定 QP 模式工作。如果配置为 1，ROI 将以 QP 偏移模式工作。\regindex{x} 可以是 A 或 B，表示对应于视频序列 A 或 B。
    \item \hyperref[regdesc:H264AROIREGION0REG]{H264\_\regindex{x}\_ROI\_REGION\regindex{n}\_REG}：此寄存器存储每个矩形 ROI 的启用和位置信息。\regindex{x} 可以是 A 或 B，表示对应于视频序列 A 或 B。\regindex{n} 可以是 0 到 7，表示 8 个 ROI。从 ROI 0 到 ROI 7，优先级依次降低。当存在重叠区域时，优先级最高的 ROI 的配置生效。此寄存器包含以下字段：
        \begin{itemize}
            \item \hyperref[fielddesc:H264AROIREGION0EN]{H264\_\regindex{x}\_ROI\_REGION\regindex{n}\_EN}：设置为 1 以启用 ROI \regindex{n}，设置为 0 以关闭 ROI \regindex{n}。当寄存器 \hyperref[fielddesc:H264AROIEN]{H264\_\regindex{x}\_ROI\_EN} 设置为 1 时，此寄存器可以控制 ROI \regindex{n} 的开关，否则此寄存器无效。
            \item \hyperref[fielddesc:H264AROIREGION0X]{H264\_\regindex{x}\_ROI\_REGION\regindex{n}\_X}：配置一个视频帧中 ROI \regindex{n} 的水平起始宏块坐标。
            \item \hyperref[fielddesc:H264AROIREGION0Y]{H264\_\regindex{x}\_ROI\_REGION\regindex{n}\_Y}：配置一个视频帧中 ROI \regindex{n} 的垂直起始宏块坐标。
            \item \hyperref[fielddesc:H264AROIREGION0XLEN]{H264\_\regindex{x}\_ROI\_REGION\regindex{n}\_X\_LEN}：配置一个视频帧中 ROI \regindex{n} 的水平宏块数量。
            \item \hyperref[fielddesc:H264AROIREGION0YLEN]{H264\_\regindex{x}\_ROI\_REGION\regindex{n}\_Y\_LEN}：配置一个视频帧中 ROI \regindex{n} 的垂直宏块数量。
        \end{itemize}
    \item \hyperref[fielddesc:H264AROIREGION0QP]{H264\_\regindex{x}\_ROI\_REGION\regindex{n}\_QP}：此寄存器用于配置每个矩形 ROI 中宏块的 QP 调整量。\regindex{x} 可以是 A 或 B，表示对应于视频序列 A 或 B。\regindex{n} 可以是 0 到 7，分别指代 8 个 ROI 之一。此寄存器的配置方法与 \hyperref[fielddesc:H264AROIMODE]{H264\_\regindex{x}\_ROI\_MODE} 寄存器的配置值有关：
        \begin{itemize}
            \item 当 \hyperref[fielddesc:H264AROIMODE]{H264\_\regindex{x}\_ROI\_MODE} 设置为 0 时：表示 ROI 模式为固定 QP 模式。此时，\hyperref[fielddesc:H264AROIREGION0QP]{H264\_\regindex{x}\_ROI\_REGION\regindex{n}\_QP} 寄存器的值直接作为对应 ROI 中宏块的 QP 值，范围为 10 到 51。
            \item 当 \hyperref[fielddesc:H264AROIMODE]{H264\_\regindex{x}\_ROI\_MODE} 设置为 1 时：表示 ROI 模式为 QP 偏移模式。此时，\hyperref[fielddesc:H264AROIREGION0QP]{H264\_\regindex{x}\_ROI\_REGION\regindex{n}\_QP} 寄存器的第 0 位到第 5 位表示 QP 偏移，范围为 0 到 51。第 6 位表示偏移的方向，0 表示正偏移，1 表示负偏移。
        \end{itemize}
    \item \hyperref[fielddesc:H264ANOROIREGIONQP]{H264\_\regindex{x}\_NO\_ROI\_REGION\_QP}：此寄存器用于配置 ROI 以外宏块的 QP 偏移。第 0 位到第 5 位表示 QP 偏移量，范围为 0 到 51。第 6 位表示偏移的方向，0 表示正偏移，1 表示负偏移。\regindex{x} 可以是 A 或 B，表示对应于视频序列 A 或 B。当 ROI 功能被禁用时，即 \hyperref[fielddesc:H264AROIEN]{H264\_\regindex{x}\_ROI\_EN} 寄存器设置为 0 时，此寄存器不生效，ROI 以外宏块的 QP 不会改变。
\end{itemize}

注意，任何 ROI 调整后的 QP 值都不应小于 10。

无论宏块是在 ROI 中还是 ROI 以外，其最终调整的 QP 值都限制在 0 到 51 之间，以确保不超出 H264 标准范围。


\subsection{H264 专用 DMA}\label{sec:h264-dma}
H264 专用 DMA 能够根据软件配置，控制编码器算法核与内存之间的数据传输。

\subsubsection{架构}\label{sec:h264-dma-arc}
%%%%
H264\_DMA 的架构如图 \ref{fig:h264_dma_architecture} 所示：

\begin{figure}[H]
    \centering
    \includegraphics[width=1.0\textwidth]{\modulefiles/figures/h264_dma_architecture.png}
    \caption{H264\_DMA 架构}
    \label{fig:h264_dma_architecture}
\end{figure}

每个模块的功能如下：
\begin{itemize}
    \item register：选择和配置寄存器。
    \item in\_dscr：向内部存储器发起一个 AXI 读取请求，获取 RX 通道的链表描述符，并将描述符发送到相应 RX 通道的 in\_link 模块。
    \item in\_link：接收对应通道的 in\_dscr 模块中获取的描述符，并将其拆分为多个 AXI 写数据请求，以引导相应的 in\_fifo 模块将数据写入内部或外部存储器。RX通道 0/1/4 用于将数据写入外部存储器，RX 通道 2/3/5 用于将数据写入内部存储器。
    \item in\_fifo：存储由 h264 或 TX 通道 3/4 发送的数据，并等待 AXI 主机读取相应通道的数据。RX 通道 0/1/2/3/4 的 in\_fifo 存储由 h264 发送的数据，而 RX 通道 5 的 in\_fifo 存储由 TX 通道 3/4 发送的数据。
    \item out\_dscr：向内部存储器发起一个 AXI 读取请求，获取 TX 通道的链表描述符，并将描述符发送到相应 TX 通道的 out\_link 模块。
    \item out\_link：接收对应通道的 out\_dscr 模块中获取的描述符，并将其拆分为多个 AXI 读数据请求，来从内部或外部存储器获取数据。TX 通道 0/3/4 从外部存储器读取数据，TX 通道 1/2 从内部存储器读取数据。
    \item out\_fifo：接收 AXI 读请求返回的数据，并等待 h264 或 RX 通道 5 读取。TX 通道 0/1/2 用于 h264 读取，TX 通道 3/4 用于 RX 通道 5 读取。
    \item arb：接收对外部存储器或内部存储器的读/写请求，并对这些读写请求进行仲裁。经过仲裁后，将相关请求分别发送到外部存储器的 rdcmd\_fifo 和 wrcmd\_fifo 模块，以及内部存储器的 rdcmd\_fifo 和 wrcmd\_fifo 模块。
    \item wrcmd\_fifo：缓存由 arb 模块仲裁的 AXI 写请求，并将写请求和数据切割成多个 8 字节数据块，然后发送到 async\_in\_fifo 模块。
    \item rdcmd\_fifo：缓存由 arb 模块仲裁的 AXI 读请求，并将读请求发送到 async\_in\_fifo 模块。
    \item bak\_fifo：缓存由 arb 模块仲裁的读/写请求信息，包括通道 id、AXI\_id、读/写请求数据长度，读/写请求请求源和读请求起始地址。该模块根据 async\_out\_fifo 模块返回的 AXI\_id 找到请求的对应信息，并将请求的完成信息和 async\_out\_fifo 模块中读取的数据发送到相应通道。
    \item async\_in\_fifo：接收来自 wrcmd\_fifo 模块和 rdcmd\_fifo 模块的 AXI 请求，并将 H264\_DMA 内核时钟域的请求同步到 AXI 时钟域，等待 AXI 主机读取。
    \item async\_out\_fifo：接收由 AXI\_master\_inter/AXI\_master\_exter 模块返回的读数据和读/写请求的完成信息，并将其从 AXI 时钟域同步到 H264\_DMA 内核的时钟域。
    \item AXI\_master\_inter/AXI\_master\_exter：从 async\_in\_fifo 模块读取读/写请求，并将其转换为 AXI 总线协议格式输出。接收由 AXI 从机返回的读数据和读/写完成信息，并将其转换为内部私有接口格式，写入 async\_out\_fifo 模块。AXI\_master\_inter 模块用于访问内部存储空间，AXI\_master\_exter 模块用于访问外部存储空间。

\end{itemize}


\subsubsection{链表描述符}\label{sec:h264-dma-link}
%%%%
H264\_DMA 链表描述符如图 \ref{fig:h264_dma_descriptor} 所示：

\begin{figure}[H]
    \centering
    \includegraphics[width=0.6\textwidth]{\modulefiles/figures/h264_dma_descriptor_p4.png}
    \caption{H264\_DMA 链表描述符}
    \label{fig:h264_dma_descriptor}
\end{figure}

\begin{itemize}
\item owner：表示与当前描述符对应的缓冲区允许的操作者。
    \begin{itemize}
    \item 0：允许操作者为 CPU。
    \item 1：允许操作者为 H264\_DMA。H264\_DMA 在使用描述符后，会清除此位。
    \item 对发送描述符，由于发送描述符默认使用后不会写回，所以默认的 owner 位不会被清除。如果需要此描述符的 owner 位在使用后自动清除，需配置 H264\_DMA\_OUT\_AUTO\_WRBACK 启用写回功能。
    \end{itemize}

\item eof：结束标志。
    \begin{itemize}
    \item 0：当前描述符不是最后一个描述符。
    \item 1：当前描述符是最后一个描述符。
    \end{itemize}
\item mod：数据块读取模式。
    \begin{itemize}
    \item 0：H264\_DMA 仅读取/写入一个水平宽度为 hb、垂直高度为 vb 的图像块一次。
    \item 1：H264\_DMA 读取/写入水平宽度为 hb、垂直高度为 vb 的图像块多次。直到读/写完具有 HA 水平宽度和 VA 垂直高度的图像数据。注意，hb/vb/HA/VA的单位是像素。
    \item 在 1D 模式下，此字段必须为 0。
    \end{itemize}

\item err\_eof：接收结束错误标志。
    \begin{itemize}
    \item 对于接收描述符，硬件在接收到一个数据包并检测到接收数据错误后，会将此位设置为 1。
    \item 对于发送描述符，此位固定为 0。
    \end{itemize}

\item 2DEN：2D 功能使能信号。
    \begin{itemize}
    \item 0：H264\_DMA 的功能与一般 DMA 功能相同。此时，可以用来传输连续的数据段。
    \item 1：启用 H264\_DMA 功能，并根据配置的图像块进行数据传输。
    \end{itemize}

\item length[13:0]/HB：
    \begin{itemize}
    \item 2DEN 为 0 时，表示传输数据的长度（最低 14 位）。对于数据接收，此字段在硬件接收到一帧数据后填充。对于数据传输，此字段表示要发送的数据长度，单位为字节。
    \item 2DEN 为 1 时，表示图像块的水平宽度 (HB)，单位为像素。
    \end{itemize}

\item size[13:0]/VB：
    \begin{itemize}
    \item 2DEN 为 0 时，表示为传输数据分配的内存空间大小（最低 14 位），单位为字节。
    \item 2DEN 为 1 时，表示图像块的垂直高度 (VB)，单位是像素。
    \end{itemize}
\item length[28:14]/HA：
    \begin{itemize}
    \item 2DEN 为 0 时，表示传输数据的长度（最高 14 位，只使用 14 \verb+~+ 27 位）。对于数据接收，此字段在硬件接收到一帧数据后填充。对于数据传输，此字段表示要发送的数据长度，单位是字节。
    \item 2DEN 为 1 时，表示描述符对应图像的水平宽度 (HA)，单位是像素。
    \end{itemize}
\item size[27:14]/VA：
    \begin{itemize}
    \item 2DEN 为 0 时，表示传输数据的内存空间大小（最高 14 位），单位是字节。
    \item 2DEN 为 1 时，表示描述符对应图像的垂直高度 (VA)，单位是像素。
    \end{itemize}
\item Buffer address pointer：
    \begin{itemize}
    \item 2DEN 为 0 时，表示 buffer 地址指针。
    \item 2DEN 为 1 时，表示 buffer 中 HA x VA 图像块原点坐标的地址指针。
    \end{itemize}
\item Next descriptor address：下一个描述符地址的指针。
\end{itemize}


\subsubsection{传输启动}\label{sec:h264-dma-start}
%%%%

H264\_DMA 根据与通道链表描述符相关的寄存器配置，从内部存储器读取链表到本地，使得相应的通道可以用这些链表传输数据，从而启动相应的通道。

与链表描述符相关的寄存器包括：
\begin{itemize}
    \item H264\_DMA\_OUTLINK/INLINK\_START\_CH\regindex{x}：此寄存器控制 H264\_DMA 的相应通道是否开始获取描述符。将此寄存器设置为 1，表示相应的通道开始获取描述符。获取描述符后，硬件自动将此寄存器设置为 0，然后硬件自动获取后续的描述符。\regindex{x} 可以为 0、1、2、3、4，表示对应通道 0、1、2、3、4 的寄存器。
    \item H264\_DMA\_OUTLINK/INLINK\_ADDR\_CH\regindex{x}：此寄存器表示存储第一个描述符的地址，需要 8 字节对齐。

\end{itemize}

\subsubsection{传输复位}\label{sec:h264-dma-reset}
%%%%
为确保 H264\_DMA 安全复位，防止通道卡住，H264\_DMA 通道的复位须按照以下步骤进行：
\begin{enumerate}
    \item 向相应通道的 H264\_DMA\_OUT/IN\_CMD\_DISABLE\_CH\regindex{x} 寄存器写入 1。
    \item 等待相应通道的 H264\_DMA\_OUT/IN\_RESET\_AVAIL\_CH\regindex{x} 寄存器置位。
    \item 向相应通道的 H264\_DMA\_OUT/IN\_RST\_CH\regindex{x} 寄存器写入 1。
    \item 向相应通道的 H264\_DMA\_OUT/IN\_RST\_CH\regindex{x} 寄存器写入 0。
\end{enumerate}

与通道复位相关的寄存器如下：
\begin{itemize}
    \item H264\_DMA\_OUT/IN\_CMD\_DISABLE\_CH\regindex{x}：此寄存器控制 H264\_DMA 的相应通道是否停止向 AXI 发起读/写请求。将此寄存器设置为 1，表示相应通道停止向 AXI 发起读/写请求。将此寄存器设置为 0，表示相应通道可以向 AXI 发起读/写请求。\regindex{x} 在 TX 通道中可以是 0、1 和 2，表示对应于通道 0、1 和 2。在 RX 通道中，\regindex{x} 可以是 0、1、2、3、4 和 5，表示对应通道 0、1、2、3、4 和 5。
    \item H264\_DMA\_OUT/IN\_RESET\_AVAIL\_CH\regindex{x}：此寄存器指示 H264\_DMA 的相应通道是否可以安全复位。如果此寄存器为 1，表示相应的通道可以安全复位。如果此寄存器为 0，表示相应的通道无法安全复位。
    \item H264\_DMA\_OUT/IN\_RST\_CH\regindex{x}：此寄存器控制 H264\_DMA 的相应通道是否执行软复位。将此寄存器设置为 1，表示相应的通道执行软复位。将此寄存器设置为 0，表示相应的通道不执行软复位。
\end{itemize}

注意：
\begin{itemize}
    \item 由于 TX 通道 3、TX 通道 4 和 RX 通道 5 互相绑定，这三个通道只有一组 H264\_DMA\_OUT/IN\_CMD\_DISABLE\_CH5、H264\_DMA\_OUT/IN\_RESET\_AVAIL\_CH5 和 H264\_DMA\_OUT/IN\_RST\_CH5 寄存器。在复位这三个通道时，只需操作这一组寄存器，即可对这三个 H264\_DMA 通道进行安全复位。
    \item 为确保数据传输的完成，在复位 RX 通道 0、RX 通道 1、RX 通道 3 和 RX 通道 4 前，需要检查相应的通道链表描述符是否写回成功。只有在通道的链表描述符写回成功后，即通道的所有数据传输完成后，才能复位通道。但是，当未启用 MV 合并功能，或 GOP 中的视频帧数为 1 时（即寄存器 H264\_GOP\_NUM 配置为 1 时），无需检查 RX 通道 3 的链表描述符是否写回成功。
\end{itemize}

\subsubsection{通道配置}\label{sec:h264-dma-ch-cfg}
%%%%
在软件启动 H264\_DMA 前，需要先对 H264\_DMA 进行配置，包括寄存器配置和链表描述符配置。

链表描述符配置：需要配置描述符的通道包括 TX 通道 0 \verb+~+ TX 通道 4，RX 通道 0 \verb+~+ RX 通道 4。链表描述符均为循环链表。
\begin{enumerate}
    \item 配置读取原始图像的通道（TX 通道 0）：
    \begin{itemize}
        \item 模式：2D，mode1。
        \item 长度/大小：HA 代表图像的水平像素数，VA 代表图像的垂直像素数，HB 为 16 x 4，VB 为 16。
    \end{itemize}
    \item 配置读取参考图像的通道（TX 通道 1）：
    \begin{itemize}
        \item 模式：2D，mode1。
        \item 长度/大小：HA 为 NUM\_H\_MB x 16 x 16，VA 为 3，HB 为 16 x 16，VB 为 3。
    \end{itemize}
    \item 配置读取 DB 中间数据的通道（TX 通道 2）：
    \begin{itemize}
        \item 模式：1D。
        \item 长度/大小：读取空间大小为 NUM\_H\_MB x (16 x 4 + 8 x 4 x 2)。
    \end{itemize}
    \item 配置读取 DB 数据前 12 行的通道（TX 通道 3）：
    \begin{itemize}
        \item 模式：1D。
        \item 长度/大小：读取空间大小为 (NUM\_H\_MB x (图像中的垂直宏块行数 - 1) x (16 x 12 + 8 x 4 x 2)) + (NUM\_H\_MB x 16 x 16 x 1.5)。
    \end{itemize}
    \item 配置读取 DB 数据的最后 4 行的通道（TX 通道 4）：
    \begin{itemize}
        \item 模式：1D。
        \item 长度/大小：读取空间大小是 (NUM\_H\_MB x (图像中的垂直宏块行数 - 1) x (16 x 4 + 8 x 4 x 2))。
    \end{itemize}
    \item 配置写入 DB 数据前 12 行的通道（RX 通道 0）：
    \begin{itemize}
        \item 模式：1D。
        \item 长度/大小：写入空间的大小与 TX 通道 3 的大小相同。
    \end{itemize}
    \item 配置写入 DB 数据的最后 4 行的通道（RX 通道 1）：
    \begin{itemize}
        \item 模式：1D。
        \item 长度/大小：写入空间的大小与 TX 通道 4 的大小相同。
    \end{itemize}
    \item 配置写入 DB 中间数据的通道的通道（RX 通道 2）：
    \begin{itemize}
        \item 模式：1D。
        \item 长度/大小：写入空间的大小与 TX 通道 2 的大小相同。
    \end{itemize}
    \item 配置写入 MV 合并数据的通道（RX 通道 3）：
    \begin{itemize}
        \item 模式：1D。
        \item 长度/大小：写入空间大小是 NUM\_H\_MB x 垂直 宏块 行数 x 4。
    \end{itemize}
    \item 配置写入比特流数据的通道（RX 通道 4）：
    \begin{itemize}
        \item 模式：1D。
        \item 长度/大小：写入空间的大小是图像的水平像素数 x 垂直像素行数 x 1.5。
    \end{itemize}
\end{enumerate}

寄存器配置：
\begin{itemize}
    \item H264\_DMA\_OUT/IN\_MEM\_BURST\_LENGTH\_CH\regindex{x}：此寄存器控制相应通道的 AXI 突发长度。
    \begin{itemize}
        \item 0：8 字节。1：16 字节。2：32 字节。3：64 字节。4：128 字节。
        \item 在 TX 通道中，\regindex{x} 可以是 0 \verb+~+ 4。在 RX 通道中，\regindex{x} 可以是 0 \verb+~+ 5。
        \item TX 通道 2 和 RX 通道 2 的突发长度配置必须相同。
    \end{itemize}
    \item \hyperref[fielddesc:H264DMAOUTREORDERENCH0]{H264\_DMA\_OUT\_REORDER\_EN\_CH0}：此寄存器控制 TX 通道 0 是否启用数据重排序。当 TX 通道 0 的链表描述符配置 HB 不为 16 时，必须将此寄存器配置为 1。
    \item \hyperref[fielddesc:H264DMABLOCKSTARTADDRCH5]{H264\_DMA\_BLOCK\_START\_ADDR\_CH5}：此寄存器控制 RX 通道 5 的数据起始地址。此寄存器配置必须与 TX 通道 1 链表描述符中的数据缓冲区地址一致。
    \item \hyperref[fielddesc:H264DMABLOCKROWLENGTH4LINECH5]{H264\_DMA\_BLOCK\_ROW\_LENGTH\_4LINE\_CH5}：此寄存器指定一行 4 行块中包含的字节数，为 (NUM\_H\_MB x (16 x 4 + 8 x 4 x 2))。
    \item \hyperref[fielddesc:H264DMABLOCKROWLENGTH12LINECH5]{H264\_DMA\_BLOCK\_ROW\_LENGTH\_12LINE\_CH5}：此寄存器指定一行 12 行块中包含的字节数，为 (NUM\_H\_MB x (16 x 12 + 8 x 4 x 2))。
    \item \hyperref[fielddesc:H264DMABLOCKLENGTH4LINECH5]{H264\_DMA\_BLOCK\_LENGTH\_4LINE\_CH5}：此寄存器指定一个 4 行块包含的字节数，为 (16 x 4 + 8 x 4 x 2)。
    \item \hyperref[fielddesc:H264DMABLOCKLENGTH12LINECH5]{H264\_DMA\_BLOCK\_LENGTH\_12LINE\_CH5}：此寄存器指定一个 12 行块包含的字节数，为 (16 x 12 + 8 x 4 x 2)。
\end{itemize}

%%%%%%%%%%%%%%%%%%
%%% Interrupts %%%
%%%%%%%%%%%%%%%%%%

\section{中断}
\label{sec:h264enc-interrupt}
\chipname{} 的 H264 编码器可以生成以下中断信号，并将其发送到\hyperref[mod:intmtrx]{中断矩阵}。
    \begin{itemize}
        \item H264\_REG\_INT
        \item H264\_DMA\_IN\_CH0\_INT
        \item H264\_DMA\_IN\_CH1\_INT
        \item H264\_DMA\_IN\_CH2\_INT
        \item H264\_DMA\_IN\_CH3\_INT
        \item H264\_DMA\_IN\_CH4\_INT
        \item H264\_DMA\_IN\_CH5\_INT
        \item H264\_DMA\_OUT\_CH0\_INT
        \item H264\_DMA\_OUT\_CH1\_INT
        \item H264\_DMA\_OUT\_CH2\_INT
        \item H264\_DMA\_OUT\_CH3\_INT
        \item H264\_DMA\_OUT\_CH4\_INT
    \end{itemize}

H264 编码器有多个内部中断源可以生成上述中断信号。

H264\_REG\_INT 中断由以下内部中断源生成：
\begin{itemize}
\item H264\_DMA\_MOVE\_2MB\_LINE\_DONE\_INT：当 H264\_DMA 将两个宏块行的参考图像从外部存储移动到内部存储时，触发此中断。
\item H264\_FRAME\_DONE\_INT：当 ENC\_CORE 完成一帧视频的压缩和编码时，触发此中断。
\item H264\_REC\_READY\_INT：当 ENC\_CORE 将 \hyperref[fielddesc:H264ARECREADYTRIGGERMBLINES]{H264\_\regindex{x}\_REC\_READY\_TRIGGER\_MB\_LINES} 行宏块的去块滤波结果（即参考图像）写入外部存储后，触发此中断。\regindex{x} 可以是 A 或 B，表示该寄存器对应于视频序列 A 或 B。
\item H264\_DB\_TMP\_READY\_INT：当 ENC\_CORE 将 \hyperref[fielddesc:H264ADBTMPREADYTRIGGERMBNUM]{H264\_\regindex{x}\_DB\_TMP\_READY\_TRIGGER\_MB\_NUM} 个宏块的去块滤波中间数据写入内部存储后，触发此中断。\regindex{x} 可以是 A 或 B，表示该寄存器对应于视频序列 A 或 B。
\end{itemize}

要触发上述中断源，必须先将中断使能寄存器 \hyperref[fielddesc:H264DMAMOVE2MBLINEDONEINTENA]{\regindex{XXX}\_ENA} 设置为 1。要清除中断，将寄存器 \hyperref[fielddesc:H264DMAMOVE2MBLINEDONEINTCLR]{\regindex{XXX}\_CLR} 设置为 1（\regindex{XXX} 为中断源名称）。

H264\_DMA\_IN\_CH\regindex{x}\_INT（\regindex{x} 为 0 \verb+~+ 5）中断由以下中断源生成：

\begin{itemize}
\item H264\_DMA\_IN\_DONE\_CH\regindex{x}\_INT：当 RX 通道 \regindex{x} (0 \verb+~+ 5) 接收到与接收描述符对应的所有数据时，触发此中断。
\item H264\_DMA\_IN\_SUC\_EOF\_CH\regindex{n}\_INT：当接收描述符中的 suc\_eof 位为 1，并且 RX 通道 \regindex{x} (0 \verb+~+ 5) 接收到与此接收描述符对应的数据时（即已接收到与 inlink 对应的数据帧或数据包），触发此中断。
\item H264\_DMA\_IN\_ERR\_EOF\_CH\regindex{x}\_INT：当 RX 通道 \regindex{x} (0 \verb+~+ 4) 接收的数据帧或数据包中检测到错误时，触发此中断。
\item H264\_DMA\_IN\_DSCR\_EMPTY\_CH\regindex{x}\_INT：当接收描述符指向的缓冲区大小小于通 RX 通道 \regindex{x} (0 \verb+~+ 4) 要接收的数据长度时，触发此中断。
\item H264\_DMA\_IN\_DSCR\_ERR\_CH\regindex{x}\_INT：当 RX 通道 \regindex{x} (0 \verb+~+ 4) 上的接收描述符未通过两项描述符检查中的任意一项时，触发此中断。
\item H264\_DMA\_INFIFO\_OVF\_L1\_CH\regindex{x}\_INT：当 RX 通道 \regindex{x} (0 \verb+~+ 5) 的 L1 RX FIFO 溢出时，触发此中断。
\item H264\_DMA\_INFIFO\_OVF\_L2\_CH\regindex{x}\_INT：当 RX 通道 \regindex{x} (0 \verb+~+ 4) 的 L2 RX FIFO 溢出时，触发此中断。
\item H264\_DMA\_INFIFO\_UDF\_L1\_CH\regindex{x}\_INT：当 RX 通道 \regindex{x} (0 \verb+~+ 5) 的 L1 RX FIFO 下溢时，触发此中断。
\item H264\_DMA\_INFIFO\_UDF\_L2\_CH\regindex{x}\_INT：当 RX 通道\regindex{x} (0 \verb+~+ 4) 的 L2 RX FIFO 下溢时,触发此中断。
\item H264\_DMA\_IN\_DSCR\_TASK\_OVF\_CH\regindex{x}\_INT：当 RX 通道\regindex{x} (0 \verb+~+ 4) 的描述符就绪任务 FIFO 溢出时，触发此中断。
\item H264\_DMA\_FETCH\_MB\_COL\_CNT\_OVF\_CH\regindex{x}\_INT：当 RX 通道 \regindex{x} (\regindex{x} = 5) 获取宏块列计数溢出时，触发此中断。
\end{itemize}

H264\_DMA\_OUT\_CH\regindex{x}\_INT（\regindex{x} 为 0 \verb+~+ 4）信号由以下中断源生成：

\begin{itemize}
\item H264\_DMA\_OUT\_DONE\_CH\regindex{x}\_INT：当与发送描述符对应的所有数据都通过 TX 通道 \regindex{x} 发送后，触发此中断。
\item H264\_DMA\_OUT\_EOF\_CH\regindex{x}\_INT：当发送描述符中的 suc\_eof 位为 1，且与此描述符对应的数据已通过 TX 通道 \regindex{x} 发送时，触发此中断。如果 H264\_DMA\_OUT\_EOF\_MODE\_CH\regindex{x} 为 0，此中断将在与此描述符对应的数据的最后一个字节进入 H264\_DMA 的 TX 通道时触发；如果 H264\_DMA\_OUT\_EOF\_MODE\_CH\regindex{x} 为 1，此中断将在从 H264\_DMA 的 TX 通道取出最后一个字节时触发。
\item H264\_DMA\_OUT\_DSCR\_ERR\_CH\regindex{x}\_INT：当 TX 通道 \regindex{x} 上的发送描述符未通过两项描述符检查中的任意一项时，触发此中断。
\item H264\_DMA\_OUT\_TOTAL\_EOF\_CH\regindex{x}\_INT：当通过 TX 通道 \regindex{x} 发送了与链表（包括多个描述符）对应的所有数据时，触发此中断。
\item H264\_DMA\_OUTFIFO\_OVF\_L1\_CH\regindex{x}\_INT：当 H264\_DMA 的 L1 TX FIFO 溢出时，触发此中断。
\item H264\_DMA\_OUTFIFO\_OVF\_L2\_CH\regindex{x}\_INT：当 H264\_DMA 的 L2 TX FIFO 溢出时，触发此中断。
\item H264\_DMA\_OUTFIFO\_UDF\_L1\_CH\regindex{x}\_INT：当 H264\_DMA的 L1 TX FIFO 下溢时，触发此中断。
\item H264\_DMA\_OUTFIFO\_UDF\_L2\_CH\regindex{x}\_INT：当 H264\_DMA 的 L2 TX FIFO 下溢时，触发此中断。
\item H264\_DMA\_OUT\_DSCR\_TASK\_OVF\_CH\regindex{x}\_INT：当 TX 通道 \regindex{x} 的描述符就绪任务 FIFO 溢出时，触发此中断。
\end{itemize}

%%%%%%%%%%%%%%%%%%%%%%%%%%%%%%
%%% Programming Procedures %%%
%%%%%%%%%%%%%%%%%%%%%%%%%%%%%%

\section{配置流程}
\label{sec:h264enc-prog-procedures}

H264 编码器具有两种工作模式：GOP 模式和双流模式。在 GOP 模式下，编码器在指定的 P 帧间隔中插入一个 I 帧。在双流模式下，编码器轮流处理两个视频序列，硬件会自动切换相应的 ENC\_CORE 配置。

\subsection{GOP 模式}\label{sec:h264-gop-cfg}
在 GOP 模式下，编码单个视频序列 A，图像的编码顺序如图 \ref{fig:h264-gop-frame-flow} 所示（每个 GOP 包含 3 个图像）。

\begin{figure}[H]
    \centering
    \includegraphics[width=1.0\textwidth]{\modulefiles/figures/h264-gop-frame-flow}
    \caption{GOP 模式下的图像编码顺序}
    \label{fig:h264-gop-frame-flow}
\end{figure}

在 GOP 模式下，前缀为 H264\_A\_ 的寄存器有效，前缀为 H264\_B\_ 的寄存器无效。

GOP 模式的配置过程如下：
\begin{enumerate}
    \item 参考章节 \ref{sec:h264-dma-ch-cfg} 配置每个 H264\_DMA 通道：
    \begin{enumerate}
        \item 配置 H264\_DMA TX 通道 0。
        \item 配置 H264\_DMA TX 通道 1。
        \item 配置 H264\_DMA TX 通道 2。
        \item 配置 H264\_DMA TX 通道 3 和 TX 通道 4。
        \item 配置 H264\_DMA RX 通道 0 和 RX 通道 1。
        \item 配置 H264\_DMA RX 通道 2。
        \item 配置 H264\_DMA RX 通道 3。
        \item 配置 H264\_DMA RX 通道 4。
        \item 配置 H264\_DMA RX 通道 5。
    \end{enumerate}
    \item\label{item:h264-gop-cfg} 根据所需功能配置 ENC\_CORE：
    \begin{enumerate}
        \item 配置 GOP 模式：将寄存器 \hyperref[fielddesc:H264FRAMEMODE]{H264\_FRAME\_MODE} 和寄存器 \hyperref[fielddesc:H264DUALSTREAMMODE]{H264\_DUAL\_STREAM\_MODE} 都设置为 0。
        \item 配置每个 GOP 包含的图像数量：将寄存器 \hyperref[fielddesc:H264GOPNUM]{H264\_GOP\_NUM} 配置为 n，则表示一个 GOP 包含 n 个图像。当此寄存器配置为 0 时，GOP 的第一帧是 I 帧，后续帧始终是 P 帧，直到编码器被重置为止。
        \item 配置图像的宏块分辨率：将寄存器 \hyperref[fielddesc:H264ASYSTOTALMBX]{H264\_A\_SYS\_TOTAL\_MB\_X} 配置为图像水平方向的宏块数量，将寄存器 \hyperref[fielddesc:H264ASYSTOTALMBY]{H264\_A\_SYS\_TOTAL\_MB\_Y} 配置为图像垂直方向的宏块数量。水平方向的宏块数量等于视频宽度除以 16 并向上取整的值，垂直方向的宏块数量等于视频高度除以 16 并向上取整的值。
        \item 配置图像初始 QP 值：将寄存器 \hyperref[fielddesc:H264AQP]{H264\_A\_QP} 配置为初始 QP 值，范围为 10 到 51。
        \item 配置去块滤波器：将寄存器 \hyperref[fielddesc:H264ABYPASSDBFILTER]{H264\_A\_BYPASS\_DB\_FILTER} 设置为 0 以启用去块滤波器，将寄存器 \hyperref[fielddesc:H264ABYPASSDBFILTER]{H264\_A\_BYPASS\_DB\_FILTER} 设置为 1 以禁用去块滤波器。
        \item 配置切片头信息：根据 H264 标准，切片头需包含一些必要信息，这些信息可在以下寄存器中进行配置：\hyperref[fielddesc:H264SLICERMEAINBIT]{H264\_SLICE\_RMEAIN\_BIT}、\hyperref[fielddesc:H264SLICEREMAINBITLENGTH]{H264\_SLICE\_REMAIN\_BITLENGTH}、\hyperref[fielddesc:H264SLICEBYTELENGTH]{H264\_SLICE\_BYTE\_LENGTH}、\hyperref[fielddesc:H264SLICEBYTELSB]{H264\_SLICE \_BYTE\_LSB}，和\hyperref[fielddesc:H264SLICEBYTEMSB]{H264\_SLICE\_BYTE\_MSB}。例如切片头是一个 41 位变量，最高有效位是切片头的第一位，那么这些寄存器应该按以下方式设置：
        \begin{itemize}
                \item \hyperref[fielddesc:H264SLICEBYTEMSB]{H264\_SLICE\_BYTE\_MSB}：切片头的第 41 位到第 10 位
                \item \hyperref[fielddesc:H264SLICEBYTELSB]{H264\_SLICE\_BYTE\_LSB}：切片头的第 9 位到第 2 位
                \item \hyperref[fielddesc:H264SLICERMEAINBIT]{H264\_SLICE\_RMEAIN\_BIT}：切片头的第 1 位
                \item \hyperref[fielddesc:H264SLICEBYTELENGTH]{H264\_SLICE\_BYTE\_LENGTH}：5
                \item \hyperref[fielddesc:H264SLICEREMAINBITLENGTH]{H264\_SLICE\_BYTE\_REMAIN\_BITLENGTH}：1
            \end{itemize}
        \item 如需使用量化结果降采样功能，请参考 \ref{sec:h264-decimate} 章节配置相关寄存器。
        \item 如需使用宏块级码率控制功能，请参考 \ref{sec:h264-rate-ctrl} 章节配置相关寄存器。
        \item 如需使用 ROI 功能，请参考 \ref{sec:h264-roi} 章节配置相关寄存器。
        \item 如需使用 MV 合并功能，请参考 \ref{sec:h264-mv-merge} 章节配置相关寄存器。
    \end{enumerate}
    \item\label{item:h264-gop-frame} 根据当前要编码的视频帧的类型，执行相应的操作。（注意，在一个 GOP 中，第一帧是 I 帧，后续帧都是 P 帧）。
    \begin{itemize}
        \item 当前待编码帧为 I 帧：
        \begin{enumerate}
            \item 启动 H264\_DMA TX 通道 0：请参考 \ref{sec:h264-dma-start} 章节进行配置。
            \item 启动 H264\_DMA RX 通道 0 和 RX 通道 1：请参考 \ref{sec:h264-dma-start} 章节进行配置。
            \item 启动 H264\_DMA RX 通道 2：请参考 \ref{sec:h264-dma-start} 章节进行配置。
            \item 如要使用 MV 合并功能，启动 H264\_DMA RX 通道 3：请参考 \ref{sec:h264-dma-start} 进行配置。
            \item 启动 H264\_DMA RX 通道 4：请参考 \ref{sec:h264-dma-start} 章节进行配置。
            \item 清除 H264 中断并启用中断。
            \item 开始图像编码：将寄存器 \hyperref[fielddesc:H264FRAMESTART]{H264\_FRAME\_START} 写为 1。
            \item 等待 264\_DB\_TMP\_READY\_INT 中断。中断触发后，启动 H264\_DMA TX 通道 2。请参考 \ref{sec:h264-dma-start} 章节进行配置。
            \item 等待 H264\_REC\_READY\_INT 中断。中断触发后，启动 H264\_DMA R 通道 5 并启动 H264\_DMA TX 通道 3 和 TX 通道 4。请参考 \ref{sec:h264-dma-start} 章节进行配置。
            \item 启动 H264\_DMA 将参考图像从外部存储移动到内部存储：将寄存器 \hyperref[fielddesc:H264DMAMOVESTART]{H264\_DMA\_MOVE\_START} 写为 1。
            \item 等待 H264\_DMA\_MOVE\_2MB\_LINE\_DONE\_INT 中断。中断触发后，启动 H264\_DMA TX 通道 1。配置方法请参考 \ref{sec:h264-dma-start} 章节的描述。
            \item 等待 H264\_FRAME\_DONE\_INT 中断。中断触发表示图像的编码已完成。
            \item 重置 H264\_DMA TX 通道 2 和 RX 通道 2：配置方法请参考 \ref{sec:h264-dma-reset} 章节的描述。
            \item 重置 264\_DMA 中解块滤波器中间数据的读写次数计数器：将寄存器 \hyperref[fielddesc:H264DMARXCH2INTERCOUNTERRST]{H264\_DMA\_RX\_CH2\_INTER\_COUNTER\_RST} 设置为 1，然后设置为 0。
        \end{enumerate}
        \item 当前待编码帧为 P 帧：
        \begin{enumerate}
            \item 清除 H264 中断并启用中断。
            \item 启动 H264\_DMA RX 通道 2：配置方法请参考 \ref{sec:h264-dma-start} 章节的描述。
            \item 开始图像编码：将寄存器 \hyperref[fielddesc:H264FRAMESTART]{H264\_FRAME\_START} 写为 1。
            \item 等待 H264\_DB\_TMP\_READY\_INT 中断。中断触发时，启动 H264\_DMA TX 通道 2。配置方法请参考 \ref{sec:h264-dma-start} 章节的描述。
            \item 等待 H264\_FRAME\_DONE\_INT 中断。中断触发时，表示图像编码已完成。
            \item 重置 H264\_DMA TX 通道 2 和 RX 通道 2：配置方法请参考 \ref{sec:h264-dma-reset} 章节的描述。
            \item 重置 H264\_DMA 中去块滤波器中间数据的读写次数计数器：将寄存器 \hyperref[fielddesc:H264DMARXCH2INTERCOUNTERRST]{H264\_\-DMA\_\-RX\_\-CH2\_\-INTER\_\-COUNTER\_\-RST} 设置为 1，然后设置为 0。
        \end{enumerate}
        \item 可以在任何帧编码结束后重新配置以下寄存器：
        \begin{itemize}
            \item 与量化结果降采样功能相关的寄存器，参考 \ref{sec:h264-decimate} 章节的描述。
            \item 与宏块级别码率控制功能相关的寄存器，参考\ref{sec:h264-rate-ctrl}章节的描述。
            \item 去块滤波器使能寄存器：\hyperref[fielddesc:H264ABYPASSDBFILTER]{H264\_A\_BYPASS\_DB\_FILTER}。
            \item ROI 功能相关的寄存器，参考 \ref{sec:h264-roi} 章节的描述。
            \item 与切片头信息相关的寄存器：\hyperref[fielddesc:H264SLICERMEAINBIT]{H264\_SLICE\_RMEAIN\_BIT}、\hyperref[fielddesc:H264SLICEREMAINBITLENGTH]{H264\_SLICE\_REMAIN\_BITLENGTH}、\hyperref[fielddesc:H264SLICEBYTELENGTH]{H264\_SLICE\_BYTE\_LENGTH}、\hyperref[fielddesc:H264SLICEBYTELSB]{H264\_SLICE\_BYTE\_LSB} 和 \hyperref[fielddesc:H264SLICEBYTEMSB]{H264\_SLICE\_BYTE\_MSB}。
            \item 与 MV 合并功能相关的寄存器，参考 \ref{sec:h264-mv-merge} 章节的描述。
        \end{itemize}
    \end{itemize}
    \item 重复步骤 \ref{item:h264-gop-frame}，直到一个 GOP 中的所有视频帧都完成编码。
    \item 确保数据完整写入。方法请参考 \ref{sec:h264-dma-reset} 章节的描述。
    \item 重置 H264\_DMA 的所有发送通道和接收通道，参考 \ref{sec:h264-dma-reset} 章节的描述。
    \item 重置 H264\_DMA 中去块滤波器结果数据的读写次数计数器，以及参考图像数据的读写次数计数器。具体而言，将寄存器 \hyperref[fielddesc:H264DMARXCH0EXTERCOUNTERRST]{H264\_DMA\_RX\_CH0\_EXTER\_COUNTER\_RST} 设置为 1，然后设置为 0；将寄存器 \hyperref[fielddesc:H264DMARXCH1EXTERCOUNTERRST]{H264\_\-DMA\_\-RX\_\-CH1\_\-EXTER\_\-COUNTER\_\-RST} 设置为 1，然后设置为 0；将寄存器\hyperref[fielddesc:H264DMARXCH5INTERCOUNTERRST]{H264\_\-DMA\_\-RX\_\-CH5\_\-INTER\_\-COUNTER\_\-RST} 设置为 1，然后设置为 0。
    \item 从步骤 \ref{item:h264-gop-cfg} 开始，重新配置并编码下一个 GOP。
\end{enumerate}


\subsection{双流模式}\label{sec:h264-dual-cfg}
在双流模式下，编码两个视频序列 A 和 B，图像的编码顺序如图 \ref{fig:h264-dual-frame-flow} 所示（一个 GOP 包含 3 个图像）。

\begin{figure}[H]
    \centering
    \includegraphics[width=1.0\textwidth]{\modulefiles/figures/h264-dual-frame-flow}
    \caption{双流模式下图像编码的顺序}
    \label{fig:h264-dual-frame-flow}
\end{figure}

在双流模式下，H264 编码帧的顺序固定，为 A1，B1，A2，B2，A3，B3，…，依此类推。在处理视频序列 A 的帧时，自动应用前缀为 H264\_A\_ 的寄存器，在处理视频序列 B 的帧时，自动应用前缀为 H264\_B\_ 的寄存器。其他前缀的寄存器对两个视频序列的帧都有效。

双流模式的配置过程如下：
\begin{enumerate}
    \item\label{item:h264-dual-cfg} 根据所需功能，配置 ENC\_CORE：
    \begin{enumerate}
        \item 配置双流模式：将寄存器 \hyperref[fielddesc:H264FRAMEMODE]{H264\_FRAME\_MODE} 设置为 1，并将寄存器 \hyperref[fielddesc:H264DUALSTREAMMODE]{H264\_DUAL\_STREAM\_MODE} 设置为 1。
        \item 配置 GOP（视频序列 A 和 B 的 GOP 配置必须相同）：将寄存器H264\_GOP\_NUM 配置为 n，表示一个 GOP 包含 n 个图像。当寄存器配置为 0 时，GOP 的第一帧是 I 帧，后续帧始终是 P 帧，直到编码器复位为止。
        \item 配置图像的宏块分辨率：将寄存器 \hyperref[fielddesc:H264ASYSTOTALMBX]{H264\_A\_SYS\_TOTAL\_MB\_X} 配置为视频序列 A 在水平方向的宏块数，将寄存器 \hyperref[fielddesc:H264ASYSTOTALMBY]{H264\_A\_SYS\_TOTAL\_MB\_Y} 配置为视频序列 A 在垂直方向的宏块数。将寄存器 \hyperref[fielddesc:H264BSYSTOTALMBX]{H264\_B\_SYS\_TOTAL\_MB\_X} 配置为视频序列 B 在水平方向的宏块数，将寄存器 \hyperref[fielddesc:H264BSYSTOTALMBY]{H264\_B\_SYS\_TOTAL\_MB\_Y} 配置为视频序列 B 在垂直方向的宏块数。水平方向的宏块数等于视频宽度除以 16 并向上取整，垂直方向的宏块数等于视频高度除以 16 并向上取整。
        \item 配置图像的初始 QP：将寄存器 \hyperref[fielddesc:H264AQP]{H264\_A\_QP} 配置为视频序列 A 的初始 QP，范围为 10 到 51。将寄存器 \hyperref[fielddesc:H264BQP]{H264\_B\_QP} 配置为视频序列 B 的初始 QP，范围为 10 到 51。
        \item 配置去块滤波器：将寄存器 \hyperref[fielddesc:H264ABYPASSDBFILTER]{H264\_A\_BYPASS\_DB\_FILTER} 设置为 0，以启用视频序列 A 的去块滤波器。将寄存器 \hyperref[fielddesc:H264ABYPASSDBFILTER]{H264\_A\_BYPASS\_DB\_FILTER} 设置为 1，以禁用视频序列 A 的去块滤波器。将寄存器 \hyperref[fielddesc:H264BBYPASSDBFILTER]{H264\_B\_BYPASS\_DB\_FILTER} 设置为 0，以启用视频序列 B 的去块滤波器。将寄存器 \hyperref[fielddesc:H264BBYPASSDBFILTER]{H264\_B\_BYPASS\_DB\_FILTER} 设置为 1，以禁用视频序列 B 的去块滤波器。
        \item 配置切片头信息：根据 H264 标准，切片头需包含一些必要信息，配置这些信息的相关寄存器包括：\hyperref[fielddesc:H264SLICERMEAINBIT]{H264\_SLICE\_RMEAIN\_BIT}，\hyperref[fielddesc:H264SLICEREMAINBITLENGTH]{H264\_SLICE\_REMAIN\_BITLENGTH}、\hyperref[fielddesc:H264SLICEBYTELENGTH]{H264\_SLICE\_BYTE\_LENGTH}、\hyperref[fielddesc:H264SLICEBYTELSB]{H264\_SLICE \_BYTE\_LSB}，和\hyperref[fielddesc:H264SLICEBYTEMSB]{H264\_SLICE\_BYTE\_MSB}。例如，如果切片头是一个 41 位变量，最高有效位是切片头的第一位，那么这些寄存器应按以下方式设置：
            \begin{itemize}
                \item \hyperref[fielddesc:H264SLICEBYTEMSB]{H264\_SLICE\_BYTE\_MSB}：切片头的第 41 位到第 10 位
                \item \hyperref[fielddesc:H264SLICEBYTELSB]{H264\_SLICE\_BYTE\_LSB}：切片头的第 9 位到第 2 位
                \item \hyperref[fielddesc:H264SLICERMEAINBIT]{H264\_SLICE\_RMEAIN\_BIT}：切片头的第 1 位
                \item \hyperref[fielddesc:H264SLICEBYTELENGTH]{H264\_SLICE\_BYTE\_LENGTH}：5
                \item \hyperref[fielddesc:H264SLICEREMAINBITLENGTH]{H264\_SLICE\_BYTE\_REMAIN\_BITLENGTH}：1
            \end{itemize}
        \item 如需使用量化结果降采样功能，请参考 \ref{sec:h264-decimate} 章节配置相关寄存器。
        \item 如需使用配置宏块级码率控制功能，请参考 \ref{sec:h264-rate-ctrl} 章节配置相关寄存器。
        \item 如需使用 ROI 功能，请参考 \ref{sec:h264-roi} 章节配置相关寄存器。
        \item 如需使用 MV 合并功能，请参考 \ref{sec:h264-mv-merge} 章节配置相关寄存器。
    \end{enumerate}
    \item\label{item:h264-dual-dma} 根据待编码的帧所属的视频序列（A 或 B），使用与视频序列对应的信息，并参考 \ref{sec:h264-dma-ch-cfg} 章节的描述配置每个 H264\_DMA 通道：
    \begin{enumerate}
        \item 配置 H264\_DMA TX 通道 0。
        \item 配置 H264\_DMA TX 通道 1。
        \item 配置 H264\_DMA TX 通道 2。
        \item 配置 H264\_DMA TX 通道 3 和 TX 通道 4。
        \item 配置 H264\_DMA RX 通道 0 和 RX 通道 1。
        \item 配置 H264\_DMA RX 通道 2。
        \item 配置 H264\_DMA RX 通道 3。
        \item 配置 H264\_DMA RX 通道 4。
        \item 配置 H264\_DMA RX 通道 5。
    \end{enumerate}
    \item\label{item:h264-dual-frame} 根据当前要编码的视频帧的类型，执行相应的操作。视频序列 A 和 B 交替编码，即首先编码视频序列 A 中的一帧，然后编码视频序列 B 中的一帧，交替进行。对于任何视频序列，在一个 GOP 中，第一帧是 I 帧，后续帧是 P 帧。
    \begin{itemize}
        \item 当前待编码帧为 I 帧：
        \begin{enumerate}
            \item 启动 H264\_DMA TX 通道 0：配置方法请参考 \ref{sec:h264-dma-start} 章节的描述。
            \item 启动 H264\_DMA RX 通道 0 和 RX 通道 1：配置方法请参考 \ref{sec:h264-dma-start} 章节的描述。
            \item 启动 H264\_DMA RX 通道 2：配置方法请参考 \ref{sec:h264-dma-start} 章节的描述。
            \item 启动 H264\_DMA RX 通道 4：配置方法请参考 \ref{sec:h264-dma-start} 章节的描述。
            \item 清除 H264 中断并启用中断。
            \item 开始图像编码：将寄存器 \hyperref[fielddesc:H264FRAMESTART]{H264\_FRAME\_START} 写为 1。
            \item 等待 H264\_DB\_TMP\_READY\_INT 中断。中断触发时，启动H264\_DMA TX 通道 2。配置方法请参考 \ref{sec:h264-dma-start} 章节的描述。
            \item 等待 H264\_FRAME\_DONE\_INT 中断。中断触发时，表示图像的编码已完成。
            \item 确保数据完全写入。方法请参考 \ref{sec:h264-dma-reset} 章节的描述。
            \item 重置 H264\_DMA 的每个发送和接收通道，参考 \ref{sec:h264-dma-reset} 章节的描述。
            \item 重置 H264\_DMA 中去块滤波器中间数据的读写计数器，重置参考图像数据的读写计数器。将寄存器 \hyperref[fielddesc:H264DMARXCH2INTERCOUNTERRST]{H264\_DMA\_RX\_CH2\_INTER\_COUNTER\_RST} 设置为 1，然后设置为 0。将寄存器 \hyperref[fielddesc:H264DMARXCH5INTERCOUNTERRST]{H264\_DMA\_RX\_CH5\_INTER\_COUNTER\_RST} 设置为 1，然后设置为 0。
        \end{enumerate}
        \item 当前待编码帧为 P 帧：
        \begin{enumerate}
            \item 清除 H264 中断并启用中断。
            \item 启动 H264\_DMA TX 通道 0：配置方法请参考 \ref{sec:h264-dma-start} 章节的描述。
            \item 启动 H264\_DMA RX 通道 0 和 RX 通道 1：配置方法请参考 \ref{sec:h264-dma-start} 章节的描述。
            \item 启动 H264\_DMA RX 通道 2：配置方法请参考 \ref{sec:h264-dma-start} 章节的描述。
            \item 启动 H264\_DMA RX 通道 4：配置方法请参考 \ref{sec:h264-dma-start} 章节的描述。
            \item 启动 H264\_DMA RX 通道 5：配置方法请参考 \ref{sec:h264-dma-start} 章节的描述。
            \item 启动 H264\_DMA TX 通道 3 和 TX 通道 4：配置方法请参考 \ref{sec:h264-dma-start} 章节的描述。
            \item 如需使用 MV 合并功能，应启动 H264\_DMA RX 通道 3：配置方法请参考 \ref{sec:h264-dma-start} 章节的描述。
            \item 启动 H264\_DMA，将参考图像从外部存储移动到内部存储：将寄存器 \hyperref[fielddesc:H264DMAMOVESTART]{H264\_DMA\_MOVE\_START} 写为 1。
            \item 等待 H264\_DMA\_MOVE\_2MB\_LINE\_DONE\_INT 中断。中断触发时，启动 H264\_DMA TX 通道 1。配置方法请参考 \ref{sec:h264-dma-start} 章节的描述。
            \item 开始图像编码：将寄存器 \hyperref[fielddesc:H264FRAMESTART]{H264\_FRAME\_START} 写为 1。
            \item 等待 H264\_DB\_TMP\_READY\_INT 中断。中断触发时，启动 H264\_DMA TX 通道 2。配置方法请参考 \ref{sec:h264-dma-start} 章节的描述。
            \item 等待 H264\_FRAME\_DONE\_INT 中断。中断触发时，表示图像编码已完成。
            \item 确保数据完全写入。方法请参考 \ref{sec:h264-dma-reset} 章节的描述。
            \item 重置 H264\_DMA 的每个发送和接收通道，参考 \ref{sec:h264-dma-reset} 章节的描述。
            \item 重置 H264\_DMA 中去块滤波器中间数据的读写计数器和参考图像数据的读写计数器。将寄存器 \hyperref[fielddesc:H264DMARXCH2INTERCOUNTERRST]{H264\_DMA\_RX\_CH2\_INTER\_COUNTER\_RST} 设置为 1，然后设置为 0。将寄存器 \hyperref[fielddesc:H264DMARXCH5INTERCOUNTERRST]{H264\_DMA\_RX\_CH5\_INTER\_COUNTER\_RST} 设置为 1，然后设置为 0。
        \end{enumerate}
        \item 在任何帧编码结束后，可以重新配置以下与当前编码的视频序列对应的寄存器：
        \begin{itemize}
            \item 与量化结果降采样函数相关的寄存器，请参考 \ref{sec:h264-decimate} 章节的描述。
            \item 与宏块级别码率控制功能相关的寄存器，请参考 \ref{sec:h264-rate-ctrl} 章节的描述。
            \item 去块滤波器使能寄存器：\hyperref[fielddesc:H264ABYPASSDBFILTER]{H264\_A\_BYPASS\_DB\_FILTER}。
            \item ROI 功能相关寄存器，请参考 \ref{sec:h264-roi} 章节的描述。
            \item 与切片头信息相关的寄存器：\hyperref[fielddesc:H264SLICERMEAINBIT]{H264\_SLICE\_RMEAIN\_BIT}、\hyperref[fielddesc:H264SLICEREMAINBITLENGTH]{H264\_SLICE\_REMAIN\_BITLENGTH}、\hyperref[fielddesc:H264SLICEBYTELENGTH]{H264\_SLICE\_BYTE\_LENGTH}、\hyperref[fielddesc:H264SLICEBYTELSB]{H264\_SLICE\_BYTE\_LSB} 和 \hyperref[fielddesc:H264SLICEBYTEMSB]{H264\_SLICE\_BYTE\_MSB}。
            \item 与 MV 合并功能相关的寄存器，请参考 \ref{sec:h264-mv-merge} 章节的描述。
        \end{itemize}
    \end{itemize}
    \item 重复步骤 \ref{item:h264-dual-dma} 和 \ref{item:h264-dual-frame}，直到两个视频序列都完成一个 GOP 编码。
    \item 重置 H264\_DMA 中去块滤波器结果数据的读写次数计数器。将寄存器 \hyperref[fielddesc:H264DMARXCH0EXTERCOUNTERRST]{H264\_DMA\_RX\_CH0\_EXTER\_COUNTER\_RST}设置为 1，然后设置为 0。将寄存器\hyperref[fielddesc:H264DMARXCH1EXTERCOUNTERRST]{H264\_DMA\_RX\_CH1\_EXTER\_COUNTER\_RST}设置为 1，然后设置为 0。
    \item 从步骤 \ref{item:h264-dual-cfg} 开始，重新配置并编码下一个 GOP。
\end{enumerate}


\subsection{软重置}\label{sec:h264-soft-rst}
无论 H264 处于何种工作模式，按照以下流程重置，即可以随时停止编码。
\begin{enumerate}
    \item 重置 ENC\_CORE：在寄存器 \hyperref[fielddesc:H264SYSRSTPULSE]{H264\_SYS\_RST\_PULSE} 中写入 1。
    \item 重置 H264\_DMA 的每个发送和接收通道，参考章节 \ref{sec:h264-dma-reset}。
    \item 重置 H264\_DMA 中的每个读写次数计数器。将寄存器 \hyperref[fielddesc:H264DMARXCH0EXTERCOUNTERRST]{H264\_DMA\_RX\_CH0\_EXTER\_COUNTER\_RST} 设置为 1，然后设置为 0。将寄存器\hyperref[fielddesc:H264DMARXCH1EXTERCOUNTERRST]{H264\_DMA\_RX\_CH1\_EXTER\_COUNTER\_RST} 设置为 1，然后设置为 0。将寄存器\hyperref[fielddesc:H264DMARXCH2INTERCOUNTERRST]{H264\_DMA\_RX\_CH2\_INTER\_COUNTER\_RST} 设置为 1，然后设置为 0。将寄存器 \hyperref[fielddesc:H264DMARXCH5INTERCOUNTERRST]{H264\_DMA\_RX\_CH5\_INTER\_COUNTER\_RST} 设置为 1，然后设置为 0。
    \item 此时，H264 编码器返回到初始状态，可以根据需要重新配置。
\end{enumerate}


%%%%%%%%%%%%%%%%%%%%%
%%% Memory Blocks %%%
%%%%%%%%%%%%%%%%%%%%%






% >>> If your module has no registers, delete sample text
%       for Sections Register Summary and Registers below <<<

%%%%%%%%%%%%%%%%%%%%%%%%
%%% Register Summary %%%
%%%%%%%%%%%%%%%%%%%%%%%%

% >>> Update the sample text below and insert
%     the info on register summary into the subfile <<<

\clearpage
\section{寄存器列表}

本节中的地址是相对于 \textbf{H264 编码器} 基地址和 \textbf{H264 DMA} 基地址的相对地址。如需了解基地址详细信息，请参考 \ref{mod:sysmem} \textit{\nameref{mod:sysmem}} 章节，表 \ref{tab:sysmem-base-address} \textit{\nameref{tab:sysmem-base-address}}。

参考章节 \hyperref[glossary-access-types]{\textit{寄存器访问类型}}，了解 \textbf{访问} 列中的缩写含义。


\hypertarget{h264enc-reg-summ}{}
\subsection{H264 编码器寄存器列表}
\label{sec:h264enc-reg-summ}


\subfile{\modulefiles/40-H264ENC-reg-summ__CN}


\hypertarget{h264enc-dma-reg-summ}{}
\subsection{H264 DMA 寄存器列表}
\label{sec:h264enc-dma-reg-summ}

\subfile{\modulefiles/40-H264ENC-DMA-reg-summ__CN}


%%%%%%%%%%%%%%%%%
%%% Registers %%%
%%%%%%%%%%%%%%%%%

% >>> Update the sample text below and insert
%      the info on registers into the subfile <<<

\clearpage
\section{寄存器}

本节中的地址是相对于 \textbf{H264 编码器} 基地址和 \textbf{H264 DMA} 基地址的相对地址。如需了解基地址详细信息，请参考 \ref{mod:sysmem} \textit{\nameref{mod:sysmem}} 章节，表 \ref{tab:sysmem-base-address} \textit{\nameref{tab:sysmem-base-address}}。

如需了解如何对保留字段进行编程，请参考章节 \hyperref[programming-reserved-field]{\textit{寄存器保留字段编程}}。


\subsection{H264 编码器寄存器}

\subfile{\modulefiles/40-H264ENC-reg__CN}


\clearpage
\subsection{H264 DMA 寄存器}

\subfile{\modulefiles/40-H264ENC-DMA-reg__CN}


\end{document}
